%% ---------------------------------------------------------------------------
% This presentation is separated by sections and subsections
\section{Introduction}

\begin{frame}
    \begin{itemize}%[<+-| alert@+>] 
    	\item This is a slide template created by latex for XJTLUers.
        \item Overleaf  \\ \url{https://www.overleaf.com/latex/templates/xjtlu-beamer-template/sfrvnnpcsmgh}
        \item GitHub \\ \url{https://github.com/yaoshanliang/XJTLU-Beamer-Template}
    \end{itemize}
    
    \note {Write your notes.\\}
    \begin{note}
        {Write your notes here}
    \end{note}
\end{frame}

\section{Literature Review} 

\subsection{Usage}

\begin{frame}{Usage} 
    \begin{itemize}
        \item Beamer is a powerful and flexible \LaTeX{} class to create great looking presentations.   \url{https://www.overleaf.com/learn/latex/Beamer}
        \item Modify from Template Beamer UFC \cite{origin}
    \end{itemize}
\end{frame}

\subsection{Features}

\begin{frame}{Features}
	The visual design follows \href{https://www.xjtlu.edu.cn/en/about/professional-services/centre-for-knowledge-and-information/university-marketing-and-communications/visual-identity-assets}{\emph{VISUAL IDENTITY ASSETS}} from XJTLU.
    \begin{enumerate}
        \item {XJTLU \example{NAVY} (RGB: 1, 54, 68)}
        \item {XJTLU \example{PURPLE} (RGB: 206, 87, 193)}
    \end{enumerate}
\end{frame}


\section{Methodology}

\subsection{Blocks}
\begin{frame}{Blocks}
    % Blocks styles
    \begin{block}{Block I}
        Text
    \end{block}

    \begin{alertblock}{Block II}
        Text
    \end{alertblock}

    \begin{exampleblock}{Block III}
        Text
    \end{exampleblock}  
    
    \successbox{Success box}
    \alertbox{Alert box}
    \simplebox{Simple box}
 
\end{frame}


%% ---------------------------------------------------------------------------
\subsection{Alorgithms}
\begin{frame}{Algorithms (pseudocode)}
    \begin{algorithm}[H]
        \SetAlgoLined
        \LinesNumbered
        \SetKwInOut{Input}{input}
        \SetKwInOut{Output}{output}
        \Input{x: float, y: float}
        \Output{r: float}
        \While{True}{
          r = x + y\;
          \eIf{r >= 30}{
           ``O valor de $r$ é maior ou iqual a 10.''\;
           break\;
           }{
           ``O valor de $r$ = '', r\;
          }
         } 
         \caption{Algorithm Example}
    \end{algorithm}
\end{frame}

\begin{frame}{Algorithms}
    \lstset{language=Python}
    \lstinputlisting[language=Python]{code/main.py}
\end{frame}

\subsection{Equations}

\begin{frame}{Equation}
    \begin{block}{Equation without numbers} 
        \begin{equation*}
            J(\theta) = \mathbb{E}_{\pi_\theta}[G_t] = \sum_{s\in\mathcal{S}} d^\pi (s)V^\pi(s)=\sum_{s\in\mathcal{S}} d^\pi(s)\sum_{a\in\mathcal{A}}\pi_\theta(a|s)Q^\pi(s,a)
        \end{equation*}
    \end{block}
%    \begin{exampleblock}{Multiple equations\footnote{If containing text in equations，use $\backslash$mathrm\{\} or $\backslash$text\{\}}}
%       
%        \begin{align}
%            Q_\mathrm{target}&=r+\gamma Q^\pi(s^\prime, \pi_\theta(s^\prime)+\epsilon)\\
%            \epsilon&\sim\mathrm{clip}(\mathcal{N}(0, \sigma), -c, c)\nonumber
%        \end{align}
%    \end{exampleblock}
\end{frame}

\begin{frame}
    \begin{block}{Equation with numbers}
        % Taken from Mathmode.tex
        \begin{multline}
            A=\lim_{n\rightarrow\infty}\Delta x\left(a^{2}+\left(a^{2}+2a\Delta x+\left(\Delta x\right)^{2}\right)\right.\label{eq:reset}\\
            +\left(a^{2}+2\cdot2a\Delta x+2^{2}\left(\Delta x\right)^{2}\right)\\
            +\left(a^{2}+2\cdot3a\Delta x+3^{2}\left(\Delta x\right)^{2}\right)\\
            +\ldots\\
            \left.+\left(a^{2}+2\cdot(n-1)a\Delta x+(n-1)^{2}\left(\Delta x\right)^{2}\right)\right)\\
            =\frac{1}{3}\left(b^{3}-a^{3}\right)
        \end{multline}
    \end{block}
\end{frame}


%\begin{frame}[fragile]{\LaTeX{} Commands}
%    \begin{exampleblock}{Commands}
%        \centering
%        \footnotesize
%        \begin{tabular}{llll}
%            \cmd{chapter} & \cmd{section} & \cmd{subsection} & \cmd{paragraph} \\
%            Chapter & Section & Subsection & Paragraph \\\hline
%            \cmd{centering} & \cmd{emph} & \cmd{verb} & \cmd{url} \\
%            Centre Align & Emphasis & Verbatim & Hyperlink \\\hline
%            \cmd{footnote} & \cmd{item} & \cmd{caption} & \cmd{includegraphics} \\
%            Foodnote & Item & Caption & FigP\&Pic \\\hline
%            \cmd{label} & \cmd{cite} & \cmd{ref} \\
%            Label & Citing & Referring\\\hline
%        \end{tabular}
%    \end{exampleblock}
%    \begin{exampleblock}{Environment Command}
%        \centering
%        \footnotesize
%        \begin{tabular}{lll}
%            \env{table} & \env{figure} & \env{equation}\\
%            Table & Figure & Equation \\\hline
%            \env{itemize} & \env{enumerate} & \env{description}\\
%            Bullets & Numbering & Description \\\hline
%        \end{tabular}
%    \end{exampleblock}
%\end{frame}


\subsection{Figures}
\begin{frame}{Figures}
    \begin{figure}
        \centering
        \caption{Logo of XJTLU.}
        \includegraphics[scale=0.5]{libs/XJTLU_logo.png}
        \source{Description of XJTLU}
        \label{fig:XJTLU_logo}
    \end{figure}
\end{frame}

\subsection{Tables}
\begin{frame}{Tables}
\begin{table}[hbt]
  \begin{tabular}{l|cc}
    1& 2& \\
    \hline
    3& 4& \\
    5& 6&
  \end{tabular}
  \caption{}
\end{table}

\end{frame}

%% ---------------------------------------------------------------------------
% This frame show an example to insert multi-columns
\subsection{Multi-columns}
\begin{frame}{Multi-columns}
    \begin{columns}{}
        \begin{column}{0.5\textwidth}
            \justify
            É possível colocar mais de uma coluna utilizando os comandos de $\backslash$begin\{column\}\{\} e $\backslash$end\{column\}
        \end{column}
        \begin{column}{0.5\textwidth}
            \justify
            Porém, o espaçamento deve ser proporcional entre as colunas para que estas colunas não entrem em coflito. O espaçamento é dado pelo segundo argumento do $\backslash$begin.
        \end{column}
    \end{columns}   
    \begin{columns}{}
        \begin{column}{0.5\textwidth}
            \justify
            É possível colocar mais de uma coluna utilizando os comandos de $\backslash$begin\{column\}\{\} e $\backslash$end\{column\}
        \end{column}
        \begin{column}{0.5\textwidth}
            \justify
            Porém, o espaçamento deve ser proporcional entre as colunas para que estas colunas não entrem em coflito. O espaçamento é dado pelo segundo argumento do $\backslash$begin.
        \end{column}
    \end{columns}     
\end{frame}


%% ---------------------------------------------------------------------------
% Reference frames
\begin{frame}[allowframebreaks]
    \frametitle{Reference}
    \printbibliography
\end{frame}

%% ---------------------------------------------------------------------------
% Final frame
\begin{frame}{}
    \centering
    \huge{\textbf{\example{Thank You !}}}
    
    
\end{frame}